\dictchar{अँ (चंद्रबिंदु / अनुनासिक)}

\bhojentry{अँकवार}{\K{\ba\bchandra\bka\bva\bmaa\bra}}{Añkwaar}{हि. स्त्री. | [ठेठ]}{\starsThree}%
{आलिंगन, छाती से लगाना, गोद।}%
{अँकवार (Añkwaar) / कोरवा}%
{Embrace, open-armed hug, lap.}
\bhojexample{दउड़ के माई के अंकवार भर लिहलसि।}{He ran and embraced his mother.}

\bhojsubentry{अँकवारना}{\K{\ba\bchandra\bka\bva\bmaa\bra\bna\bmaa}}{Añkwaarna}{हि. क्रि. स. | [ठेठ]}{\starsTwo}%
{गले लगाना, छाती से चिपटाना।}%
{अँकवार लिहल / छाती से लगावल}%
{To embrace, clasp to chest.}

\bhojsubentry{अँकवारी}{\K{\ba\bchandra\bka\bva\bmaa\bra\bmii}}{Añkwari}{हि. स्त्री. | [ठेठ]}{\starsTwo}%
{गोद, क्रोड़, अंक।}%
{कोरा (Kora) / अँकवारी}%
{Lap, open arms.}

\bhojsubentry{अँकोर}{\K{\ba\bchandra\bka\bmo\bra}}{Añkor}{हि. पुं. | [ठेठ]}{\starsTwo}%
{वर या अतिथि को दी जाने वाली भेंट।}%
{अँकोर / भेंट}%
{Gift/present for guest or groom.}

\bhojentry{अँकुरी}{\K{\ba\bchandra\bka\bmu\bra\bmii}}{Añkuri}{हि. स्त्री. | [ठेठ/कृषि]}{\starsTwo}%
{पौधे की बेल का मुड़ा भाग, अंकुर, कांटी।}%
{अँकुरी (Añkuri) / काँटा}%
{Tendril, sprout shoot, small hook.}

\bhojsubentry{अँकुरावल}{\K{\ba\bchandra\bka\bmu\bra\bmaa\bva\bla}}{Añkuraawal}{हि. क्रि. | [ठेठ/कृषि]}{\starsTwo}%
{अंकुरित होना।}%
{अँकुरावल / कल्ला निकलब}%
{To sprout, germinate.}

\bhojentry{अँखुआ}{\K{\ba\bchandra\bkha\bmu\baa}}{Añkhua}{हि. पुं. | [ठेठ]}{\starsThree}%
{अंकुर, बीज से फूटता हुआ नया पौधा।}%
{अँखुआ (Añkhua) / कल्ला}%
{Sprout, germinating seed shoot.}

\bhojsubentry{अँखुआना}{\K{\ba\bchandra\bkha\bmu\baa\bna\bmaa}}{Añkhuana}{हि. क्रि. | [ठेठ]}{\starsTwo}%
{अंकुरित होना, अंखुआ निकलना।}%
{अँखुआइल / कल्ला निकलब}%
{To sprout, germinate.}

\bhojentry{अँगलू}{\K{\ba\bchandra\bga\bla\bmuu}}{Añgaluu}{हि. पुं. | [ठेठ]}{\starsOne}%
{दर्जी का उंगली में पहनने का छल्ला।}%
{अँगलू (Añgaluu) / अंगुष्ठाना}%
{Tailor's thimble, finger guard.}

\bhojentry{अँगुरी}{\K{\ba\bchandra\bga\bmu\bra\bmii}}{Añguri}{हि. स्त्री. | [ठेठ]}{\starsThree}%
{हाथ या पैर की उंगली।}%
{अँगुरी (Añguri) / आँगुर}%
{Finger, toe.}
\bhojexample{लइका अँगुरी पकड़ के चले लागल।}{The child started walking holding a finger.}

\bhojcompound{अँगुरी काटना}{\K{\ba\bchandra\bga\bmu\bra\bmii} \ \K{\bka\bmaa\btta\bna\bmaa}}{Añguri Kaatna}{हि. मुहा. | [ठेठ]}%
{आश्चर्य या पश्चाताप में दांतों तले उंगली दबाना}{अँगुरी काटल}{To bite fingers in awe/regret}

\bhojentry{अँगूठा}{\K{\ba\bchandra\bga\bmuu\bttha\bmaa}}{AñguuTha}{हि. पुं. | तद्भव}{\starsThree}%
{हाथ-पैर की पहली बड़ी उंगली, अंगुष्ठ।}%
{अँगूठा (AñguuTha) / अंगुठा}%
{Thumb, big toe.}

\bhojsubentry{अँगूठी}{\K{\ba\bchandra\bga\bmuu\bttha\bmii}}{AñguuThii}{हि. स्त्री. | तद्भव}{\starsThree}%
{उंगली की मुद्रिका, छल्ला।}%
{अँगूठी (AñguuThii) / छाप}%
{Finger ring.}

\bhojcompound{अँगूठा दिखाना}{\K{\ba\bchandra\bga\bmuu\bttha\bmaa} \ \K{\bda\bmi\bkha\bmaa\bna\bmaa}}{AñguuTha Dikhana}{हि. मुहा. | [ठेठ]}%
{साफ मना कर देना, धोखा देना}{अँगूठा देखावल}{To refuse flatly, renege}

\bhojentry{अँगौछा}{\K{\ba\bchandra\bga\bmau\bcha\bmaa}}{Añgaucha}{हि. पुं. | [ठेठ]}{\starsThree}%
{सूती देह पोंछने का तौलिया, गमछा।}%
{अँगौछा (Añgaucha) / गमछा}%
{Traditional cotton towel.}
\bhojexample{घाम में कांधा पर अँगौछा राखि के निकलीं।}{Carry a towel on your shoulder in the sun.}

\bhojsubentry{अँगौरी}{\K{\ba\bchandra\bga\bmau\bra\bmii}}{Añgauri}{हि. स्त्री. | [ठेठ/कृषि]}{\starsTwo}%
{अनाज के रूप में दी जाने वाली मजदूरी।}%
{अँगौरी (Añgauri)}%
{Agricultural wages paid in grain.}

\bhojentry{अँगार}{\K{\ba\bchandra\bga\bmaa\bra}}{Añgaar}{हि. पुं. | तद्भव}{\starsThree}%
{दहकता कोयला या लकड़ी।}%
{अँगार (Añgaar) / दहकत कोइला}%
{Burning coal, embers.}

\bhojsubentry{अँगारिन}{\K{\ba\bchandra\bga\bmaa\bra\bmi\bna}}{Añgaarin}{हि. स्त्री. | [ठेठ]}{\starsOne}%
{क्रोधित स्त्री, चिंगारी जैसी।}%
{अँगारिन}%
{Fiery/hot-tempered woman.}

\bhojentry{अँगेइया}{\K{\ba\bchandra\bga\bme\bai\bya\bmaa}}{Añgeiya}{हि. स्त्री. | [ठेठ]}{\starsTwo}%
{भोज-भात या दावत का निमंत्रण।}%
{अँगेइया (Añgeiya) / न्योता}%
{Feast invitation.}

\bhojentry{अँगना}{\K{\ba\bchandra\bga\bna\bmaa}}{Añgna}{हि. पुं. | [ठेठ]}{\starsThree}%
{घर के भीतर का खुला चौकोर स्थान।}%
{अँगना (Añgna) / आँगना}%
{Courtyard, open yard.}
\bhojexample{माई अँगना में बइठल बिया।}{Mother is sitting in the courtyard.}

\bhojsubentry{अँगनाई}{\K{\ba\bchandra\bga\bna\bmaa\bai}}{Añgnaai}{हि. स्त्री. | [ठेठ]}{\starsTwo}%
{आंगन का विस्तार या परिसर।}%
{अँगनाई}%
{Courtyard space/premises.}

\bhojentry{अँचाना}{\K{\ba\bchandra\bca\bmaa\bna\bmaa}}{Añcaana}{हि. क्रि. स. | [ठेठ]}{\starsThree}%
{भोजन के बाद हाथ-मुंह धोना।}%
{अँचाना (Añcaana) / हाथ-मुंह धोवल}%
{To wash hands and mouth after meal.}

\bhojsubentry{अँचवना}{\K{\ba\bchandra\bca\bva\bna\bmaa}}{Añcavna}{हि. पुं. | [ठेठ]}{\starsTwo}%
{हथेली में लिया गया आचमन का पानी।}%
{अँचवना / आचमन जल}%
{Water taken in palm for rinsing mouth.}

\bhojsubentry{अँचवनी}{\K{\ba\bchandra\bca\bva\bna\bmii}}{Añcavni}{हि. स्त्री. | [ठेठ]}{\starsTwo}%
{हाथ-मुंह धोने की क्रिया।}%
{अँचवनी}%
{Rinsing action.}

\bhojentry{अँचरा}{\K{\ba\bchandra\bca\bra\bmaa}}{Añcra}{हि. पुं. | [ठेठ]}{\starsThree}%
{साड़ी का पल्लू या आंचल।}%
{अँचरा (Añcra) / पल्लू}%
{Sari pallu, skirt border.}

\bhojentry{अँझुराना}{\K{\ba\bchandra\bjha\bmu\bra\bmaa\bna\bmaa}}{Añjhuraana}{हि. क्रि. | [ठेठ]}{\starsTwo}%
{उलझना, विवाद में फंसना।}%
{अँझुराइल (Añjhuraail)}%
{To get entangled, caught in dispute.}

\bhojsubentry{अँझुराइल}{\K{\ba\bchandra\bjha\bmuu\bra\bmaa\bai\bla}}{Añjhuraail}{हि. वि. | [ठेठ]}{\starsTwo}%
{उलझा हुआ, विवादित।}%
{अँझुराइल}%
{Entangled, complicated.}

\bhojsubentry{अँझुरा}{\K{\ba\bchandra\bjha\bmu\bra\bmaa}}{Añjhura}{हि. पुं. | [ठेठ]}{\starsTwo}%
{उलझा हुआ सूत या धागा।}%
{उलझल धागा}%
{Tangled thread/yarn.}

\bhojentry{अँझुरी}{\K{\ba\bchandra\bjha\bmu\bra\bmii}}{Añjhuri}{हि. स्त्री. | [ठेठ]}{\starsThree}%
{दोनों हथेलियों का सम्पुट।}%
{अँझुरी (Añjhuri) / अंजलि}%
{Cupped palms.}
\bhojexample{अँझुरी में जल भर के पूजा कइलसि।}{Filled water in cupped palms to worship.}

\bhojentry{अँटइया}{\K{\ba\bchandra\btta\bai\bya\bmaa}}{Añtaiya}{हि. स्त्री. | [ठेठ]}{\starsTwo}%
{कपड़े या रस्सी की छोटी गांठ।}%
{अँटइया (Añtaiya) / गांठ}%
{Small knot, bundle.}

\bhojsubentry{अँठइयावल}{\K{\ba\bchandra\bttha\bai\bya\bmaa\bva\bla}}{Añthaiyaawal}{हि. क्रि. | [ठेठ]}{\starsTwo}%
{दर्द से ऐंठना या मरोड़ उठना।}%
{ऐंठल / पेट मरोड़ल}%
{To writhe or twist in pain.}

\bhojentry{अँटिया}{\K{\ba\bchandra\btta\bmi\bya\bmaa}}{Añtiya}{हि. स्त्री. | [ठेठ/कृषि]}{\starsThree}%
{फसल या घास की छोटी पूली।}%
{अँटिया (Añtiya) / आँटी}%
{Small crop sheaf.}

\bhojsubentry{अँटी}{\K{\ba\bchandra\btta\bmii}}{Añtii}{हि. स्त्री. | [ठेठ]}{\starsTwo}%
{सूत या धागे की छोटी लच्छी।}%
{अँटी / लच्छी}%
{Small skein/spool of thread.}

\bhojsubentry{अँटियावल}{\K{\ba\bchandra\btta\bmi\bya\bmaa\bva\bla}}{Añtiyaawal}{हि. क्रि. | [ठेठ/कृषि]}{\starsTwo}%
{घास या फसल की पूली बांधना।}%
{पूली बाँधल}%
{To bundle crops into sheaves.}

\bhojentry{अँडवा}{\K{\ba\bchandra\bdda\bva\bmaa}}{Añdwa}{हि. पुं. | [ठेठ]}{\starsTwo}%
{अंडी का बीज।}%
{अँडवा (Añdwa) / रेड़ी के बीज}%
{Castor seed.}

\bhojentry{अँडुवा}{\K{\ba\bchandra\bdda\bmu\bva\bmaa}}{Añduwa}{हि. पुं. | [ठेठ/पशु]}{\starsTwo}%
{बिना बधाई (असबध) किया गया बैल।}%
{अँडुवा बैल (Añduwa bail)}%
{Uncastrated bull.}

\bhojsubentry{अँडुवाई}{\K{\ba\bchandra\bdda\bmu\bva\baa\bai}}{Añduwaai}{हि. स्त्री. | [ठेठ]}{\starsTwo}%
{अड़ियलपन, सांड जैसा जिद्दी व्यवहार।}%
{जिद्द / अड़ियलपन}%
{Unruly stubbornness.}

\bhojentry{अँधियारा}{\K{\ba\bchandra\bdha\bmi\bya\bmaa\bra\bmaa}}{Añdhiyaara}{हि. पुं. | तद्भव}{\starsThree}%
{अंधकार, अंधेरा।}%
{अँधियारा (Añdhiyaara) / अंहार}%
{Darkness, gloom.}

\bhojsubentry{अँधियारी}{\K{\ba\bchandra\bdha\bmi\bya\bmaa\bra\bmii}}{Añdhiyaarii}{हि. स्त्री. | [ठेठ]}{\starsTwo}%
{अंधेरी रात, आंखों की पट्टी।}%
{अँधियारी}%
{Dark night, blindfold.}

\bhojsubentry{अँधार}{\K{\ba\bchandra\bdha\bmaa\bra}}{Añdhaar}{हि. पुं. | [ठेठ]}{\starsTwo}%
{घोर अंधकार।}%
{अँधार / घोर अंहार}%
{Deep darkness.}

\bhojsubentry{अँधियार}{\K{\ba\bchandra\bdha\bmi\bya\bmaa\bra}}{Añdhiyaar}{हि. पुं. | [ठेठ]}{\starsTwo}%
{गहरा अंधेरा।}%
{अँधियार}%
{Pitch darkness.}

\bhojentry{अँसुआ}{\K{\ba\bchandra\bsa\bmu\baa}}{Añsua}{हि. पुं. | [ठेठ]}{\starsThree}%
{आंसू, लोर।}%
{अँसुआ (Añsua) / लोर (Lor)}%
{Tears, weeping drop.}
\bhojexample{आँखि से अँसुआ ढुलके लागल।}{Tears began to fall from the eyes.}

\bhojsubentry{अँसुआई}{\K{\ba\bchandra\bsa\bmu\baa\bai}}{Añsuaai}{हि. वि. | [ठेठ]}{\starsTwo}%
{आंसुओं से भरी आंखें।}%
{अँसुआई आँखि}%
{Tearful, teary-eyed.}

\bhojsubentry{अँसुआवल}{\K{\ba\bchandra\bsa\bmu\bva\bmaa\bva\bla}}{Añsuwaaval}{हि. क्रि. | [ठेठ]}{\starsTwo}%
{आंखों में आंसू भर आना।}%
{आँखि अँसुवा गइल}%
{To well up with tears.}

\bhojsubentry{अँसुरा}{\K{\ba\bchandra\bsa\bmu\bra\bmaa}}{Añsura}{हि. पुं. | [ठेठ/कृषि]}{\starsTwo}%
{हंसिया की मुड़ी धार।}%
{हंसिया के धार}%
{Curved edge of sickle.}
